% --- MAGIC COMMENTS: CONFIGURAZIONE EDITOR --- (vedi main.tex per riferimenti e spiegazioni) ---
% !TEX root = ../../main.tex

% ________________________________________ IMPOSTAZIONI GLOBALI SPECIALI (FIXES) _______________________________________

% --------------------------------------- LISTATI CODICE (Pacchetto "listings") ----------------------------------------
\lstset{style=hkn-numbered}     % Imposta di default lo stile grafico per i listati di codice, sfruttando il custom
                                % style hkn-numbered (la versione con la numerazione delle righe attiva).
                                % Comodo per non dover sempre scrivere: \begin{lstlisting}[style=hkn-numbered]
                                % NOTA: si potrebbe impostare in generale di default per il progetto aggiungendo questa
                                % istruzione nella custom class.
                                % AGGIORNAMENTO: è stato aggiunta alla custom class l'impostazione \lstset{style=hkn},
                                % con il merge del 2026-04-16: io ho impostato la versione numerata.

% ---------------------------- SPAZIATURA ELENCHI PUNTATI E NUMERATI (Pacchetto "enumitem") ----------------------------
% \setlist{noitemsep}   % Rimuove la spaziatura verticale introdotta di default tra gli elementi degli elenchi (itemize,
                        % enumerate, description)
                        % Eccessivo, simile ad usare un valore negativo per itemsep. Usata soluzione a seguire.
\setlist{itemsep=0pt}   % Riduce (ma non rimuove del tutto) lo spazio verticale tra gli item o elementi di tutti gli
                        % elenchi a 0pt, per un aspetto più compatto. Dovrebbe essere adatto alla maggior parte delle
                        % liste (si possono modificare in locale quelle da allargare) Si applica a TUTTE le liste
                        % (itemize, enumerate, description) Es per applicare solo alle liste puntate (itemize):
                        % \setlist[itemize]{itemsep=...}

% ---------------------------------- UNICODE CHARACTER MAP: DIZIONARIO E ATTIVAZIONE -----------------------------------
% Disattivate in quanto sono state attivate in "Packages.tex": li è fornita una lunga spiegazione su queste opzioni.
% Sarebbero dovuto essere attive in questa posizione, ma sono li per questioni didattiche sul codice.
% \input{glyphtounicode.tex}          % [EXTENSION] Carica il dizionario Unicode standard ufficiale.
% \pdfgentounicode=1                  % [OVERRIDE] Interruttore che attiva la CMAP NATIVA nel PDF.

% ---------------------------------------- COLORI LINK SETUP PACCHETTO hyperref ----------------------------------------
% NOTA: setup già definito in _preamble.tex. Qui lo SOVRASCRIVO per personalizzarlo.

% IMPOSTAZIONI GLOBALI E MACRO AUSILIARIE
% NOTA: ricorda che i colori originali di LaTeX, come "blue" sono "molto sparati". Quelli importati con xcolor,
% che iniziano con la maiuscola, sono invece più tenui e adatti alla lettura.
% NOTA: RefMatt sta per Reference Matter: ossia i capitoli di reference tecnica: ToC, LoT, LoF, Index, Bib...
% Tuttavia non necessariamente a tutti questi applichiamo la correzione colori. Sono di interesse principiale ToC, LoF,
% LoT. Per ora Index non è stato inserito. Mentre in Bib ci sono solo i backref links, che è giusto siano colorati.

\newcommand{\RefMattDefaultLinksColor}{%
    \hypersetup{%
        linkcolor=blue,     % Colore per riferimenti interni (es: capitoli, figure, tabelle, ecc.).
        %                   % ATTENZIONE: colora anche l'intero TOC/LOF, essendo le righe di questi dei link.
        citecolor=blue,     % Colore per le citazioni bibliografiche.
    }%
}

\newcommand{\RefMattBlackLinksColor}{%
    \hypersetup{%
        linkcolor=black,    % Colore per riferimenti interni (es: capitoli, figure, tabelle, ecc.).
        %                   % ATTENZIONE: colora anche l'intero TOC/LOF, essendo le righe di questi dei link.
        citecolor=black,    % Colore per le citazioni bibliografiche.
    }%
}

% AVVIO SOVRASCRITTURA IMPOSTAZIONI GLOBALI
\RefMattDefaultLinksColor   % NOTA: ho usato questo approccio per scrivere i colori una volta sola. Avevo creato macro
                            % ausiliarie solo con le opzioni di interesse, ma dovevo comunque riscrivere i colori più
                            % volte.
\hypersetup{%
    colorlinks=true,        % Attiva il colore del testo per i link; rimuove i riquadri colorati (antiestetici nel PDF).
    urlcolor=blue,          % Colore per i link web esterni.
    %                       % NOTA: usare lo stesso blu dei link interni riduce il feedback visivo.
    linktoc=all             % Rende cliccabili sia il titolo che la pagina negli indici.
    %                       % ATTENZIONE:Colora tutta la riga se linkcolor è attivo.
}%

% PATCH: SPEGNIMENTO COLORI LINK NEGLI INDICI:

\begin{comment} COME PATCHARE NEL MODO GIUSTO? LE VARIE POSSIBILITÀ ANALIZZATE.

1. SISTEMA DI HOOK NATIVO DI LaTeX3: utilizzando i comandi \AddToHook.
   ESEMPIO PER HOOK COMANDI:
   \AddToHook{cmd/tableofcontents/before}{\blacklinks}
   \AddToHook{cmd/tableofcontents/after}{\bluelinks}

   ESEMPIO PER HOOK AMBIENTI (EVNIROMENTS):
   \AddToHook{env/theindex/begin}{\blacklinks}  % NON servono before e after per ambienti, perchè costituiscono gruppi
                                                % di codice chiuso di per sé.

   PROBLEMA: questo sistema nativo di hook, non funziona con tutti i comandi, in particolare non funziona con macro
   complesse, che hanno particolarità interne di precedenza e attese per chiamate o letture. Uno di questi casi è
   proprio il comando \tableofcontents, che è un wrapper complesso che chiama internamente altri comandi e file.
   NOTA: in generale funzionano molto bene con macro semplici, e la cosa utile è che non modifica o sovrascrive le macro
   originali. Il sistema di hook è proprio interno e funzionante su binari a parte.

2. PATCHING CON PACKAGE ETOOOLBOX - UTILIZZANDO I COMANDI \pretocmd e \apptocmd

   Questi comandi sono molto più robusti dei precedenti, perché funzionano

\newcommand{\ColorPatchError}[2]{%                              % Comando per stampare messaggio di errore utile.
    \ClassError{HKNdocument}{%
        Color patching #2 #1 did not work.\MessageBreak
        Check the code to understand!
    }{%
        See the GlobalSettings file, "hyperref setup" section.
    }%
}


% % 1. Indice Generale (TOC)
% \xpretocmd{\tableofcontents}{\RefsBlackLinks}{}{\ColorPatchError{ToC}{before}}
% \AddToHook{cmd/tableofcontents/after}{\RefsDefaultLinks}

% % 2. Elenco delle Figure (LOF)
% \xpretocmd{\listoffigures}{\RefsBlackLinks}{}{\ColorPatchError{LoF}{before}}
% \AddToHook{cmd/listoffigures/after}{\RefsDefaultLinks}

% % 3. Elenco delle Tabelle (LOT)
% \xpretocmd{\listoftables}{\RefsBlackLinks}{}{\ColorPatchError{LoT}{before}}
% \AddToHook{cmd/listoftables/after}{\RefsDefaultLinks}

% % 4. Indice Analitico (Index)
% % Se usi imakeidx, il comando standard è \printindex.
% % Come alternativa c'è l'ambiente theindex, che è quello nativo di LaTeX.
% % \xpretocmd{\printindex}{\RefsBlackLinks}{}{\ColorPatchError{Index}{before}}
% % \xapptocmd{\printindex}{\RefsDefaultLinks}{}{\ColorPatchError{Index}{after}}
% % \AtBeginEnvironment{theindex}{\RefsBlackLinks}     % NON servono before e after per ambienti, perchè costituiscono gruppi
% %                                                 % di codice chiuso di per sé.

% % 5. Bibliografia
% \xpretocmd{\printbibliography}{\RefsBlackLinks}{}{\ColorPatchError{Bib}{before}}
% \AddToHook{cmd/printbibliography/after}{\RefsDefaultLinks}

\end{comment}

\begin{comment} INTERFACCIA DA AGGIUNGERE ALLA CUSTOM CLASS
% ______________________________________________________________________________ INIZIO MODIFICA TIZIANO - hyperref
% ORIGINARIAMENTE IN _preamble.tex
% SPOSTATO SU CUSTOM CLASS PER CENTRALIZZARE GESTIONE HYPERREF E COLORI
% --------------------------
% HYPERREF FOR INTERNAL AND EXTERNAL LINKS
% --------------------------
% Create clickable links in PDF and set link colors
% \hypersetup{
%     colorlinks=true,
%     linkcolor=black,
%     citecolor=black,
%     urlcolor=blue
% }

% MACRO AUSILIARIE
% NOTA: ricorda che i colori originali di LaTeX, come "blue" sono "molto sparati". Quelli importati con xcolor,
% che iniziano con la maiuscola, sono invece più tenui e adatti alla lettura.

% RefMatt sta per Reference Matter: ossia i capitoli di reference tecnica: ToC, LoT, LoF, Index, Bib...

\newcommand{\RefMattDefaultLinksColor}{%
    \hypersetup{%
        linkcolor=blue,     % Colore per riferimenti interni (es: capitoli, figure, tabelle, ecc.).
        %                   % ATTENZIONE: colora anche l'intero TOC/LOF, essendo le righe di questi dei link.
    }%
}

\newcommand{\RefMattBlackLinksColor}{%
    \hypersetup{%
        linkcolor=black,    % Colore per riferimenti interni (es: capitoli, figure, tabelle, ecc.).
        %                   % ATTENZIONE: colora anche l'intero TOC/LOF, essendo le righe di questi dei link.
    }%
}

% IMPOSTAZIONI GLOBALI
\RefMattDefaultLinksColor   % NOTA: ho usato questo approccio per scrivere i colori una volta sola. Avevo creato macro
                            % ausiliarie solo con le opzioni di interesse, ma dovevo comunque riscrivere i colori più
                            % volte.
\hypersetup{%
    colorlinks=true,        % Attiva il colore del testo per i link; rimuove i riquadri colorati (antiestetici nel PDF).
    urlcolor=blue,          % Colore per i link web esterni.
    %                       % NOTA: usare lo stesso blu dei link interni riduce il feedback visivo.
    linktoc=all             % Rende cliccabili sia il titolo che la pagina negli indici.
    %                       % ATTENZIONE:Colora tutta la riga se linkcolor è attivo.
}%

% INTERFACCIA DI PERSONALIZZAZIONE PER GLI UTENTI FINALI
% --- 1. VARIABILI DI STATO (Boolean flags)
% In accordo con le precedenti versioni della custom class, il ToC è sempre stampato.
\newif\ifHKN@printLoF       % Flag: se true, stampa LoF
\newif\ifHKN@printLoT       % Flag: se true, stampa LoT
\newif\ifHKN@blackColors    % Flag: se true, usa colore nero per link in ToC, LoF, LoT.
% --- 2. DICHIARAZIONE INTERFACCIA (pgfkeys) - CON GESTIONE ERRORI

\pgfkeys{
    /HKNrefs/.is family,             % Dichiara una nuova famiglia di chiavi, nella folder /HKNrefs (refs: sta per references)
    /HKNrefs,                        % Attiva il percorso suddetto come default (solo per le prossime dichiarazioni)
    %
    % STAMPA LOF (.is choice crea un menu chiuso)
    PrintLoF/.is choice,                                    % Dichiara la chiave PrintLoF come scelta multipla (choice)
    PrintLoF/true/.code  = {\HKN@printLoFtrue},             % Se usata setta true il flag per stmapre LoF
    PrintLoF/false/.code = {\HKN@printLoFfalse},            % Se usata setta false il flag per stmapre LoF
    PrintLoF/.unknown/.code = {%                            % Gestione Eccezioni: avvisa l'utente se usa valori non ammessi
        \ClassError{HKNdocument}{%                          % Tipo di errore: errore di classe, custom class HKNdocument
            Invalid value for key PrintLoF.\MessageBreak    % Messaggio stampato nell'errore
            Allowed values are: true, false%
        }{%
            Check your \string\setupRefs\space options for typos.%   % Messaggio di help (scrivi h nel terminale
            %                                                        % dopo blocco compilazione per visualizzarlo)
        }%
    },
    PrintLoF = true,                                        % Imposta il valore di default della key: se l'utente non specifica nulla, usa questo
    % STAMPA LOT (.is choice crea un menu chiuso)
    PrintLoT/.is choice,                                    % Dichiara la chiave PrintLoT come scelta multipla (choice)
    PrintLoT/true/.code  = {\HKN@printLoTtrue},             % Se usata setta true il flag per stmapre LoT
    PrintLoT/false/.code = {\HKN@printLoTfalse},            % Se usata setta false il flag per stmapre LoT
    PrintLoT/.unknown/.code = {%                            % Gestione Eccezioni: avvisa l'utente se usa valori non ammessi
        \ClassError{HKNdocument}{%                          % Tipo di errore: errore di classe, custom class HKNdocument
            Invalid value for key PrintLoT.\MessageBreak    % Messaggio stampato nell'errore
            Allowed values are: true, false%
        }{%
            Check your \string\setupRefs\space options for typos.%   % Messaggio di help (scrivi h nel terminale
            %                                                        % dopo blocco compilazione per visualizzarlo)
        }%
    },
    PrintLoT = true,                                        % Imposta il valore di default della key: se l'utente non specifica nulla, usa questo
    % GESTIONE COLORI (.is choice crea un menu chiuso)
    BlackRefsColors/.is choice,                             % Dichiara la chiave BlackRefsColors come scelta multipla (choice)
    BlackRefsColors/true/.code  = {\HKN@blackColorstrue},   % Se usata setta true il flag per attivare il nero
    BlackRefsColors/false/.code = {\HKN@blackColorsfalse},  % Se usata setta false il flag per disattivare il nero
    BlackRefsColors/.unknown/.code = {%                     % Gestione Eccezioni: avvisa l'utente se usa valori non ammessi
        \ClassError{HKNdocument}{%                          % Tipo di errore: errore di classe, custom class HKNdocument
            Invalid value for key BlackRefsColors.\MessageBreak         % Messaggio stampato nell'errore
            Allowed values are: true, false%
        }{%
            Check your \string\setupRefs\space options for typos.%      % Messaggio di help (scrivi h nel terminale
            %                                                           % dopo blocco compilazione per visualizzarlo)
        }%
    },
    BlackRefsColors = true,                                % Imposta il valore di default della key: se l'utente non specifica nulla, usa questo
    % GESTIONE ERRORI PER CHIAVI INESISTENTI
    % Se l'utente sbaglia a digitare il NOME della chiave,scatta l'unknown generale della famiglia.
    % \pgfkeyscurrentname contiene il nome della chiave sbagliata.
    .unknown/.code = {%                                                             % Gestione Eccezioni: avvisa l'utente se usa key non ammesse
        \ClassError{HKNdocument}{%                                                  % Tipo di errore: errore di classe, custom class HKNdocument
            Unknown key '\pgfkeyscurrentname' used in Refs setup.\MessageBreak      % Messaggio stampato dall'errore
            Allowed keys are: HeaderCase, ChapterLabel, NumberLabel%
        }{%
            Please read the class documentation or check for typos.%                % Messaggio di help (scrivi h nel terminale dopo blocco compilazione per visualizzarlo)
        }%
    },
}
% ______________________________________________________________________________ FINE MODIFICA TIZIANO - hyperref
\end{comment}

% ----------------------------------------------------------------------------------------------------------------------

\begin{comment}
==============================================================================
ARCHITETTURA DELLE OPZIONI PER LA CUSTOM CLASS HKNdocument E ALTRE SPIEGAZIONI
==============================================================================

________________________________________________________________________________

                   PER CLASS VERSIONE: MERGE DEL 2026-04-16
________________________________________________________________________________

--- OPZIONI DELLA CUSTOM CLASS ---
La classe HKNdocument accetta parametri opzionali passati tra parentesi quadre nel comando iniziale:
\documentclass[opzione1, opzione2, ...]{HKNdocument}

L'ORDINE NON CONTA: Poiché la classe dichiara esplicitamente ogni singola opzione (tramite \DeclareOption), il motore
LaTeX si limita a cercare le chiavi riconosciute all'interno dell'elenco fornito: possono quindi essere inserite
nell'ordine preferito.

Le Opzioni disponibili sono:

LINGUA DEL DOCUMENTO (Imposta il pacchetto babel e le diciture interne)
- italian            % (Valore di Default)
- english

DIMENSIONE DEL FONT (Passata alla classe base 'book')
- 10pt
- 11pt
- 12pt               % (Valore di Default)

PROFONDITÀ DELL'INDICE (ToC - Table of Contents)
- toc=chapters       % (Livello 0: Mostra solo i capitoli)
- toc=sections       % (Livello 1: Mostra capitoli e sezioni)
- toc=subsections    % (Livello 2: Default, mostra fino alle sottosezioni)
- toc=subsubsections % (Livello 3: Mostra fino alle sotto-sottosezioni)

ESEMPIO DI UTILIZZO:
\documentclass[english, 10pt, toc=chapters]{HKNdocument}

--- OSSERVAZIONE IMPORTANTE SULLE OPZIONI DELLA CLASSE: OPZIONI DI BASE NON SOVRASCRIVIBILI ---
Attualmente, se l'utente inserisse opzioni standard di LaTeX nel comando \documentclass (es. openright), che non sono
riportate tra le opzioni esplicitamente dichiarate per la custom class, queste verrebbero ignorate dal compilatore
generando un warning ("Unused global option(s)").
Questo succede perché:
1. Il comando \ProcessOptions usato nella custom class scarta di default tutto ciò che non è stato esplicitamente
   definito con un \DeclareOption all'interno della classe.
2. La classe base "book", sulla quale quella personalizzata HKNdocument viene generata, viene caricata con parametri
   fissi (hardcoded) tramite \LoadClass[..., a4paper, openany]{book}. Questo garantisce un'impaginazione uniforme per
   tutti i documenti dell'associazione.

--- COME SBLOCCARE LA SOVRASCRIBILITÀ DELLE OPZIONI DI BASE ---
Se in futuro si volesse permettere all'utente di passare opzioni native della classe "book", si deve
modificare il file .cls in 3 step:
1. Aggiungere il comando a seguire, subito prima di \ProcessOptions\relax. Questo consente di catturare
    tutte le opzioni non esplicitamente dichiarate e inoltrarle alla classe 'book':
    \DeclareOption*{\PassOptionsToClass{\CurrentOption}{book}}
2. Impostare le opzioni da passare di default alla custom class, che l'utente può sovrascrivere:
    \PassOptionsToClass{a4paper,openany}{book}
3. Rimuovere i vincoli rigidi (come a4paper e openany) da \LoadClass (così possono essere sovrascritti
    alla chiamata della classe). Conviene invece lasciare il vincolo rigido sulla dimensione del font base.

Quindi in conclusione l'aggiunta sarebbe, sostituendo il blocco finale delle opzioni:
\DeclareOption*{\PassOptionsToClass{\CurrentOption}{book}} % Accetta opzioni extra per l'utente
\PassOptionsToClass{a4paper,openany}{book}                 % Default (che l'utente può sovrascrivere)
\ProcessOptions\relax                                      % Processa tutte le opzioni lette
\LoadClass[\HK@classfontsize]{book}                        % \HK@classfontsize contiene il font base.
                                                            % Ok lasciarlo qui, per sicurezza.

--- NOTA SU \HK@classfontsize ---
OPZIONE ESPLICITA HARDCODED: la dimensione base del font viene iniettata tramite la macro interna \HK@classfontsize,
garantendo sia flessibilità che sicurezza strutturale. Infatti, nonostante la macro sia passata rigidamente (hardcoded)
a \LoadClass, essa rappresenta ancora prima un'opzione esplicita della custom class, quindi può essere aggiornata
dall'utente alla chiamata della class (es. \documentclass[10pt]{HKNdocument}).

LIMITE TECNICO GENERALE: In qualsiasi documento LaTeX, la dimensione base del font non può essere ridichiarata o
modificata all'interno del documento (main.tex). Si tratta infatti di un'informazione geometrica vitale che deve essere
fornita una volta per tutte alla classe 'book' nell'esatto istante in cui si inizializza. Se potesse cambiare a
posteriori, l'impaginazione verrebbe disastrosamente compromessa. Per questo l'incorporazione hardcoded tramite macro in
\LoadClass rappresenta una Good Practice.

PROTEZIONE DELLA CUSTOM MACRO TRAMITE @: L'uso del simbolo '@' nel nome della macro \HK@classfontsize costituisce un
vincolo contro il tentativo di \renewcommand su di essa. In LaTeX, la chiocciola è considerata un carattere valido per
comporre i nomi dei comandi SOLO nei file strutturali (.cls o .sty). Se l'utente provasse a manipolare la macro nei file
.tex, il compilatore restituirebbe un errore fatale ("Undefined control sequence"), perché la chiocciola non viene
riconosciuta come parte del comando.

SBLOCCO PROTEZIONE TRAMITE @: Questo sistema si può sbloccare temporaneamente tra i comandi \makeatletter e
\makeatother, che trasformano rispettivamente la chiocciola @ ("at") in un carattere normale (consentendo la chiamata),
e poi di nuovo in un carattere speciale (ripristinando la protezione).
Ad esempio il seguente codice funzionerebbe per stampare la dimensione del font di base:
\makeatletter
    \HK@classfontsize
\makeatother
\end{comment}


\begin{comment}
    ====================================================================================================================
    RIASSUNTO DELLE IMPOSTAZIONI GLOBALI DELLO STILE DEL DOCUMENTO IMPOSTATE DALLA CUSTOM CLASS: GEOMETRIA, TIPOGRAFIA E
    STRUTTURA DEL TESTO
    ====================================================================================================================

    Questo lungo ma utile commento documenta tutte le personalizzazioni introdotte dalla custom class HKNdocument (file
    HKNdocument.cls + _preamble.tex) al progetto, a livello estetico, strutturale e geometrico, che definiscono
    l'aspetto visivo di tutto il testo inserito dall'utente.
    Vengono indicati in particolare:
    - i comandi di interesse con cui sono stati impostati, o si possono impostare
    - i valori esatti dei parametri che sono stati passati ai comandi


    -----------------------------------------------------------------------------
    SEZIONE 1: IMPOSTAZIONI PRINCIPALI RELATIVE AL CONTENUTO AGGIUNTO DALL'UTENTE
    -----------------------------------------------------------------------------

    --- 1. GEOMETRIA DELLA PAGINA E MARGINI ---
    -   FORMATO CARTA E APERTURA CAPITOLI: Impostati passando le opzioni 'a4paper' (carta formato A4) e 'openany' (i
        capitoli NON sono forzati a iniziare sulla pagina destra dispari, cioè possono in questo caso iniziare su
        qualsiasi pagina) al comando di base:
            \LoadClass[...,a4paper,openany]{book}, all'interno della custom class.
        Opzioni hardcoded, non modificabili dall'utente, per garantire uniformità estetica in tutti i documenti
        dell'associazione.
    -   MARGINI DELLA PAGINA: Fissati simmetricamente a 3 cm su tutti i lati. L'impostazione è definita caricando il
        pacchetto geometry all'interno della custom class, con l'opzione dedicata:
        \RequirePackage[a4paper,margin=3cm]{geometry}

    --- 2. TIPOGRAFIA DI BASE E COLLEGAMENTI ---
    -   DIMENSIONE DEL FONT: Il default interno è 12pt (passato a \LoadClass). L'utente può sovrascriverlo richiamando
        la classe nel main con le opzioni
        dedicate: \documentclass[10pt]{HKNdocument} (oppure 11pt, 12pt).
    -   FAMIGLIA TIPOGRAFICA: Il font di default (Computer Modern) è interamente sostituito dal font "Latin Modern"
        richiamando \RequirePackage{lmodern}.
    -   COLORI DEI COLLEGAMENTI E METADATI PDF: Gestiti dal pacchetto 'hyperref' tramite l'istruzione \hypersetup{...}.
        Nella nuova versione, la classe sfrutta questo comando anche per iniettare automaticamente nel file PDF le
        proprietà di sistema (Titolo, Autore, Soggetto). Per i link, il comando esatto applicato nel codice è:
        \hypersetup{colorlinks=true, linkcolor=black, citecolor=black, urlcolor=blue}.
        In precisione:
            * 'colorlinks=true' rimuove i riquadri colorati attorno ai link, piuttosto antiestetici, colorando
               direttamente il testo;
            * 'linkcolor' e 'citecolor' rendono neri i riferimenti interni (figure, formule) e la bibliografia
            * 'urlcolor' evidenzia in blu i collegamenti a siti web esterni.
    -   PARENTESI (DELIMITATORI MATEMATICI): I comandi nativi \left( e \right) introducono tipograficamente una
        spaziatura bianca orizzontale "extra" attorno alle parentesi scalabili, allontanandole fastidiosamente dal
        contenuto. Il comando globale \mleftright, del pacchetto "mleftright", risolve questa anomalia, sopprimendo lo
        spazio extra e rendendo le formule compatte e professionali. In precisione il pacchetto importa i comandi \mleft
        e \mright, che costituiscono le versioni corrette (per il problema di spaziatura) di \left e \right, e il
        comando \mleftright che se usato, dal punto in cui è messo in poi, si comporta come un wrapper che ridefinisce
        (usa \def, non \renewcommand essendo comandi di basso livello) i comandi \left e \right, sostituendoli di fatto
        con le versioni corrette \mleft e \mright.
        I vecchi comandi originali vengono salvati in altre macro, \mleftright@left e \mleftright@right, che essendo
        protette dalla chiocciola possono essere chiamate solo tra i comandi \makeatletter e \makeatother.
        Se si volesse ripristinare il comportamento originale, si può richiamare il comando \mleftright@restore, che
        ripristina i comandi \left e \right originali.

    --- 3. PARAGRAFI: SPAZIATURA, INDENTAZIONE E INTERLINEA ---
    Lo stile utilizzato dalla custom class HKNdocument si può descrivere in generale come "Block Style Paragraph": È uno
    stile di impaginazione (tipico del web e dei manuali tecnici o scientifici) in cui i paragrafi sono formattati come
    "blocchi" separati da spazi: invece di segnalare l'inizio di un nuovo paragrafo facendo rientrare la prima riga
    verso destra, si usa una riga vuota verticale per separarli nettamente, favorendo la leggibilità e la scansione
    visiva. Per essere implementato, sfrutta i due elementi a seguire:
    -   RIENTRO DI RIGA PER PARAGRAFI (\parindent): È lo spazio orizzontale vuoto all'inizio della prima riga di un
        nuovo paragrafo. Caricando il pacchetto 'parskip', questo valore viene forzatamente azzerato (\parindent=0pt).
    -   SPAZIATURA PARAGRAFI (\parskip): È lo spazio verticale inserito per separare due paragrafi distinti. Il
        pacchetto 'parskip' lo imposta automaticamente a un valore proporzionale alla dimensione del font, permettendo
        di separare visivamente i blocchi di testo ora che il rientro iniziale è stato rimosso. Precisazione tecnica: Il
        pacchetto imposta \parskip a un valore base di circa mezza riga di testo (0.5\baselineskip, cioè è come se si
        avesse usato: \setlength{\parskip}{0.5\baselineskip}, stessa sintassi con cui si può personalizzare. Nota che un
        numero adiacente a un registro di lunghezza è interpretato da LaTeX come un prodotto). A questo valore aggiunge
        un margine elastico (chiamato "glue" in TeX): permette al motore di dilatare o restringere impercettibilmente
        quello spazio per far arrivare l'ultima riga della pagina perfettamente a filo col margine inferiore.
    -   NOTA ULTERIORE - INTERLINEA (LINE SPACING): la custom class non modifica in alcun modo l'interlinea, affidandosi
        all'algoritmo nativo di LaTeX (Single Spacing, con fattore moltiplicativo esatto di 1.0) ottimizzato per il font
        'lmodern'. Se si volesse modificare l'interlinea, si potrebbero usare:
            * \linespread{<fattore moltiplicativo>} per modificare globalmente l'interlinea.
              Sconsigliato perché agisce di forza bruta su tutto, incluse didascalie, footnotes, ecc.
            * pacchetto "setspace": lo standard migliore, permette di gestire l'interlinea del corpo del testo, senza
              toccare il resto degli elementi.

    --- 4. FORMATTAZIONE DIDASCALIE, SOTTOLINEATURE, E LISTATI DI CODICE ---
    -   DIDASCALIE (CAPTIONS):
        Hanno un margine orizzontale forzato di 1 cm per lato, impostato tramite il comando \captionsetup{margin=1cm}
        (pacchetto "caption").
    -   BLOCCHI DI CODICE (LISTINGS): Configurati tramite \lstdefinestyle{hkn}{...}
        Ora impostato globalmente di default tramite \lstset{style=hkn}. Usano font monospazio (\ttfamily), dimensione
        \small, e i nuovi colori della palette Mu Nu Chapter, con margini rientrati di 30pt (xleftmargin=30pt,
        xrightmargin=30pt).
    -   SOTTOLINEATURE (Pacchetti 'ulem' e 'contour'):
        La profondità della sottolineatura è forzata a 1.8pt (\ULdepth). Il pacchetto 'contour' genera un "contorno di
        pulizia" (\contourlength{0.8pt}): si tratta di una brevissima interruzione della sottolineatura (o minuscolo
        alone bianco invisibile) generata attorno alle lettere discendenti (come la 'p', la 'g' o la 'q'). Ciò evita che
        la linea di sottolineatura tagli fisicamente il carattere, garantendo una leggibilità eccellente.

    --- 5. STRUTTURA E NUMERAZIONE DEGLI AMBIENTI ---
    -   PROFONDITÀ INDICE (ToC): Impostata dalla variabile \HK@tocdepth a 2 (mostra fino alle \subsection).
        È modificabile dall'utente richiamando la classe con l'opportuno parametro, essendo un'opzione aggiunta
        esplicitamente: \documentclass[toc=chapters]{HKNdocument}.
    -   AUTOMAZIONE TITOLI E COPERTINA: Il frontespizio non è più basato su comandi singoli sparsi. Se non forzato con
        \title{}, il titolo si autogenera fondendo i dati inseriti nel nuovo blocco \courseinfo{} (codice, materia).
    -   SUBORDINAZIONE DELLA NUMERAZIONE:
            * Figure, Tabelle ed Equazioni sono legate ai capitoli (es. Figura 1.1) tramite il comando
              \numberwithin{...}{chapter}.
            * Custom Box Colorati dichiarati con \DeclareHKNBox (teorema, definizione, ecc.) sono subordinati
              nativamente al capitolo sfruttando l'opzione \newtcolorbox[auto counter, number within=chapter] del
              pacchetto 'tcolorbox'.

    ------------------------------------------
    SEZIONE 2: ALTRE IMPOSTAZIONI E TECNICISMI
    ------------------------------------------
    Altre impostazioni più particolari, e questioni tecniche più complesse.

    --- 1. INTESTAZIONI (HEADER): \headheight E \topmargin ---
    Altezza del box o riga di intestazione fissata manualmente a 30pt (\headheight=30pt) per ospitare il logo piccolo di
    HKN senza generare errori "Overfull", ossia per farlo stare dentro la riga dell'intestazione in modo legittimo.
    Di conseguenza Il margine superiore (\topmargin) è compensato con \addtolength{\topmargin}{-20pt}, che diminuisce il
    margine superiore per riportarlo simmetricamente a 3 cm, come gli altri margini, dopo aver lasciato lo spazio per
    l'intestazione. In conclusione grazie a questa modifica, il testo vero e proprio inizia circa (esattamente se si
    fosse scritto -18pt, essendo di default \headheight=12pt) dove sarebbe iniziato con il margine superiore di 3cm e
    con intestazione di dimensione standard.

    --- 2. PIÈ DI PAGINA (FOOTER) - NUMERI DI PAGINA PER MAINMATTER (CONTENUTO STANDARD: vedi a seguire frontmatter) ---
    Il Footer (o piè di pagina) è l'area fissa posizionata appena sopra il margine inferiore del foglio, strutturalmente
    separata dal corpo del testo, dedicata a ospitare le informazioni di servizio, come numeri di pagina.
    Nel gergo comune, nel footer vanno anche le "note a piè di pagina", ma LaTeX gestisce codeste in un'area dedicata,
    separata dal footer, chiamata \footins (o "footnote area"), separata sia dal corpo del testo che dal footer.
    La custom class applica \fancyfoot[C]{\thepage}: Il pacchetto 'fancyhdr' assegna alla casella centrale ([C]) del
    footer il compito di stampare la macro \thepage, che contiene dinamicamente il valore del contatore della pagina
    corrente.

    --- 3. TRANSIZIONE ROMANA/ARABA E FRONTMATTER (TESTO PRELIMINARE) ---
    Il comando \frontmatter: È un interruttore strutturale dell'editoria che gestisce le pagine preliminari, ossia le
    pagine di servizio dedicate all'inserimento di testo preliminare, come ad esempio l'indice, la licenza, eccetera (in
    questo caso). Quando attivato, forza la numerazione in numeri romani minuscoli (i, ii, iii...) e sospende la
    numerazione formale dei capitoli per non intaccare il conteggio del libro principale.
    Tramite l'hook (*) \AtBeginDocument, la custom class inietta \frontmatter di nascosto prima dell'indice. Terminate
    le pagine di servizio, viene iniettato \mainmatter: questo comando ripristina i contatori, ripartendo da "1" in
    numeri arabi per il primo capitolo di testo. I comandi \frontmatter e \mainmatter sono dunque comandi che vanno
    utilizzati sempre in coppia: quello che è presente tra di essi, è il contenuto preliminare citato.

    (*) un hook è un comando che intercetta un passaggio della compilazione, e inietta in quel punto codice aggiuntivo
    per modificarne l'esito.

    --- 4. PAGINE BIANCHE E FILIGRANA DI SICUREZZA ---
    -   FRONTESPIZIO: Nella ridefinizione di \maketitle, viene chiamato il comando \thispagestyle{empty}, che "spegne"
        completamente sia l'intestazione che il piè di pagina per la singola pagina di copertina.
        Inoltre, se la classe non rileva il file "version", stampa in filigrana rossa "NON PUBBLICATO" tramite il
        pacchetto 'background'. Vedi il file CONTRIBUTING.md e 01_Packages.tex per maggiori dettagli.
    -   PAGINE VUOTE: I comandi \clearpage e \cleardoublepage sono stati sovrascritti con \renewcommand per iniettarvi
        al loro interno \thispagestyle{empty}. Questo garantisce che le pagine di riempimento (generate ad esempio per
        far iniziare i capitoli a destra) siano totalmente immacolate, prive del numero di pagina stampato.
        I comandi originali di LaTeX per la gestione delle pagine vuote sono stati conservati con il comando \let:
        \let\originalclearpage\clearpage    e    \let\originalcleardoublepage\cleardoublepage
        In questo modo si possono richiamare anche i comandi originali, se necessario.

    NOTA: GESTIONE DEI SALTI PAGINA: \clearpage, \cleardoublepage e \newpage
    In LaTeX, i comandi per andare a capo pagina hanno impatti strutturali molto diversi sulla gestione delle figure e
    delle tabelle in sospeso (flottanti, ossia non ancora stampati, e che possono essere stampati in punti diversi da
    quello in cui sono stati chiamati, a discrezione del compilatore, a meno di aver specificato diversamente con
    opzioni, come [H] per "Here", del pacchetto 'float').
        * \newpage (salto semplice):
            Interrompe la pagina corrente e sposta il testo successivo sul foglio nuovo. ATTENZIONE: NON forza la stampa
            delle immagini in sospeso. Una figura inserita prima di \newpage potrebbe tranquillamente "scivolare" nelle
            pagine successive mescolandosi al nuovo testo.
        * \clearpage (reset duro):
            Interrompe la pagina corrente MA, prima di far ripartire il testo, forza la stampa immediata di tutti gli
            elementi flottanti rimasti nella memoria di compilazione. Crea un vero e proprio blocco che impedisce alle
            immagini di invadere la sezione o il capitolo successivo.
        * \cleardoublepage (Il rispetto per la rilegatura)
            Esegue lo svuotamento totale di \clearpage, ma aggiunge un controllo geometrico per l'impaginazione
            fronte-retro: garantisce che il testo successivo inizi rigorosamente su una pagina di destra (dispari).
            Se necessario, inserisce automaticamente una pagina di sinistra vuota per ripristinare l'allineamento.

    --- 5. GESTIONE DELL'INIZIO DEI CAPITOLI (STILE "plain") ---
    Il comando \chapter è stato ridefinito (*) per iniettarvi \thispagestyle{plain}. A differenza di "empty" (passato ad
    esempio nella ridefinizone di \maketitle), lo stile 'plain' nasconde l'intestazione grafica superiore (per non
    appesantire visivamente il titolo del capitolo), ma conserva regolarmente il numero di pagina centrato in basso.

    | (*) NOTA TECNICA: RIDEFINIZIONE PROFONDA DEI CAPITOLI - HACKING DI LATEX STANDARD:
    | Capire come il comando chapter è stato ridefinito è più complesso.
    | - OSSERVAZIONE TECNICA SULL'USO DI \def AL POSTO DI \renewcommand:
    |   Si utilizza la primitiva di basso livello \def perché \renewcommand ha due limiti strutturali in questo
    |   contesto:
    |   1)  non è programmato per riconoscere o gestire le versioni con asterisco (*) dei comandi, causando errori
    |       fatali. Esistono infatti due versioni del comando: \chapter la versione standard, e \chapter* la versione
    |       "starred" (con asterisco) per i capitoli non numerati.
    |   2)  andrebbe in conflitto con la macro nativa \secdef, usata nella ridefinizione. Quest'ultima passa gli
    |       argomenti testuali crudi esigendo il "pattern matching" letterale ([#1]#2), un input grezzo che solo \def
    |       può recepire senza rompere la classe 'book'.
    | - ROUTING DEI CAPITOLI IN LATEX:
    |   Il comando \chapter funge solo da router, grazie a \secdef: a seconda che l'utente usi o meno l'asterisco,
    |   smista il lavoro a \@chapter o \@schapter.
    | - MODIFICHE APPORTATE RISPETTO AL LATEX STANDARD:
    |   1)  \chapter (Il dispatcher o router):
    |       Inietta forzatamente \thispagestyle{plain}. Rimuove l'intestazione grafica (header e logo) dalla primissima
    |       pagina di ogni capitolo per pulizia visiva, mantenendo solo il numero di pagina a piè di pagina.
    |   2)  \@chapter (Capitoli numerati normali):
    |       Inietta il comando \chaptermark{#1}. Istruisce l'header (fancyhdr) a catturare e stampare in alto
    |       esclusivamente il "Titolo Corto" (#1, passato tra parentesi quadre), mentre il titolo lungo (#2) va
    |       all'Indice, evitando che titoli troppo lunghi sbordino e si sovrappongano al logo.
    |   3)  \@schapter (Capitoli non numerati / starred chapters): Inietta il comando \chaptermark{#1} come sopra.
    |       Allo stesso tempo questo risolve un bug storico, il "bug dell'intestazione fantasma": nativamente, i
    |       capitoli con asterisco NON aggiornano l'header, il quale continuerebbe a mostrare ostinatamente il nome del
    |       capitolo precedente. Questa aggiunta forza l'header ad aggiornarsi immediatamente.

    --- 6. GESTIONE GLOBALE DEI TITOLI: CORTI VS LUNGHI (OVERRIDE STANDARD LATEX) ---
    Comportamento Standard LaTeX: Usando \chapter[Corto]{Lungo}, il "Titolo Corto" viene inviato di default sia
    all'intestazione in alto (header) che all'Indice (ToC). Comportamento Custom di questa classe:
    -   Nella ridefinizione di \@chapter, il comando \addcontentsline{toc} è stato forzato a ricevere SEMPRE il
        "Titolo Lungo" (#2), garantendo un Indice dettagliato.
    -   Il comando \chaptermark{#1}, invece, cattura e invia esclusivamente il "Titolo Corto" all'intestazione,
        mantenendola compatta ed evitando che il testo sbordi sopra il logo.

    --- 7. SPAZIATURA LISTE (Pacchetto 'enumitem' in _preamble.tex) ---
    In LaTeX nativo, le liste hanno spaziature verticali grandi di default, perché in origine (anni 70-80) erano pensate
    principalmente per separare dimostrazioni matematiche su carta stampata, che richiedevano molta "aria" visiva. Oggi,
    per liste di testo brevi, questo spazio può risultare sproporzionato. La custom class non modifica questa spaziatura
    in generale, ma solo in un caso specifico, spiegato a seguire.
    -   Il pacchetto 'enumitem' permette di stringere le liste usando tre opzioni da passare in
        \begin{<tipo_lista>}[<opzioni>]:
            * topsep : Spazio elastico tra testo normale e inizio/fine lista.
            * parsep : Spazio verticale tra due paragrafi nello stesso punto elenco.
            * itemsep: Spazio verticale vuoto tra un punto elenco e il successivo.
              NOTA: tra due item agisce itemsep + parsep.
        Esempio particolare: \begin{enumerate}[parsep=10pt, itemsep=-10pt] spazia solo i paragrafi di uno stesso item,
        ma annulla la spaziatura tra item diversi, compensando con un valore negativo l'effetto di parsep.
    -   Modifica Apportata: Il preambolo carica il pacchetto senza alterare i valori globali, ma la classe lo sfrutta
        localmente (es: nel file cclicense.tex) applicando il parametro esatto [parsep=0pt] per compattare al massimo
        il testo nelle clausole legali.


    --- 8. PALETTE COLORI PERSONALIZZATA (pacchetto "xcolor") ---
    La classe utilizza costruttori \definecolor{nome}{RGB}{R, G, B} per definire variabili globali, garantendo la
    massima coerenza estetica in tutto il progetto:

    COLORI UFFICIALI HKN:
    - hknBlue                   RGB(  0, 102, 153)
    - hknOrange                 RGB(243, 112,  33)
    - hknYellow                 RGB(255, 196,  37)
    - hknRed                    RGB(227,  27,  35)
    COLORI MU NU CHAPTER:
    - munuDarkBlue              RGB( 22,  65, 127)
    - munuMediumBlue            RGB( 48, 121, 185)
    - munuLightBlue             RGB(  0, 172, 233)
    COLORI DI SFONDO E TESTO:
    - backgroundLight           RGB(242, 242, 242)
    - textGrayMedium            RGB(146, 154, 166)

                                    COLORI VECCHI ELIMINATI CON MERGE DEL 2026-04-16
    COLORI D'ACCENTO:
    - accentYellow              RGB(254, 196,  41)
    - accentRed                 RGB(236,  45,  36)
    COLORI DI SUPPORTO:
    - supportOrange             RGB(242, 183,   5)
    - supportDarkBlue           RGB( 55,  81, 113)
    COLORI DI SFONDO:
    - backgroundLight           RGB(242, 242, 242)
    COLORI PER TESTI E BORDI (SCALA DI GRIGI/BLU):
    - textGrayBlue              RGB(100, 117, 140)
    - textGrayMedium            RGB(146, 154, 166)
    - textGrayLight             RGB(184, 187, 191)
    COLORI SPECIFICI PER I RIQUADRI TCOLORBOX:
    - tcolorboxLeftColor        RGB(  2,  65, 191)
    - tcolorboxBackTitleColor:  RGB(119, 152, 255)
    - tcolorboxBackColor        RGB(210, 226, 255)

    --- 9. IMPOSTAZIONI ESTETICHE E GRAFICHE DEI COLORED CUSTOM BOX (PACCHETTO tcolorbox) ---
    PARAMETRI GEOMETRICI E GRAFICI ESATTI:
    - boxrule=0mm           Disabilita il bordo standard su quattro lati per evitare l'effetto "gabbia" pesante.
    - leftrule=1mm          Applica una barra verticale accentuata solo sul lato sinistro, creando un richiamo visivo
                            forte al colore (parametro #4).
    - arc=0.5mm             Imposta un raggio di curvatura discreto e più spigoloso per gli angoli.
                            NOTA: prima del merge del 2026-04-16, il raggio era di 2mm, più arrotondato e morbido.
    - rounded corners=all   Forza l'arrotondamento su tutti e quattro gli angoli del box, armonizzando la struttura.
    GESTIONE CROMATICA (Identità HKN):
    - colframe=#4           Il colore del bordo sinistro e degli accenti è passato dinamicamente col parametro 4
                            (es. hknYellow) (NOTA: simile all'ex accentYellow, usato prima del merge del 2026-04-16).
    - colbacktitle          Fissato su 'textGrayMedium' per mantenere il titolo neutro e leggibile, indipendentemente
                            dal colore del box.
    - colback               Utilizza 'backgroundLight' per lo sfondo interno, garantendo contrasto massimo con il testo
                            nero.
    LOGICA FUNZIONALE:
    - [auto counter, number within=chapter]     subordina la numerazione dei box ai capitoli
                                                (es. Teorema 1.1, 1.2, 2.1, ecc.) e automatizza il conteggio.
    - \iflanguage{...}      Istruzione condizionale, rileva la lingua del documento e sceglie automaticamente tra il
                            nome italiano (parametro #2) o inglese (parametro #3) del box.
\end{comment}
